% Figure 15 revolved about the y-axis: r^2 <= y <= r, 0 <= r <= 1.
% The front half of the solid is cut away; the blue shell is retained.
% Projection: (r,theta,y) -> (2.5 r cos(theta), 3.2 y + .65 r sin(theta)).
\begin{tikzpicture}[line width=.65pt,line cap=round,line join=round,
  every node/.style={font=\small,inner sep=2pt}]
% The visible conical surface slopes down from the rim to the origin.
\fill[gray!12] (0,0)--(2.5,3.2)--
  plot[domain=0:180,samples=73]
    ({2.5*cos(\x)},{3.2+.65*sin(\x)})--cycle;
% Exposed radial sections show the curved lower boundary.
\foreach \s in {-1,1}{
  \fill[gray!24] (0,0)--({2.5*\s},3.2)--
    plot[domain=1:0,samples=45] ({2.5*\s*\x},{3.2*\x*\x})--cycle;
  \draw[gray!65] (0,0)--({2.5*\s},3.2);
  \draw[gray!65]
    plot[domain=0:1,samples=45] ({2.5*\s*\x},{3.2*\x*\x});
}
\draw[gray!65] plot[domain=0:180,samples=73]
  ({2.5*cos(\x)},{3.2+.65*sin(\x)});
\foreach \a in {45,90,135}{
  \draw[gray!40,line width=.4pt] (0,0)--
    ({2.5*cos(\a)},{3.2+.65*sin(\a)});
}
% A complete translucent radial band x <= r <= x+dx.
% Both cylindrical walls continue behind the shell as well as in front.
\foreach \r/\alpha in {.58/.10,.62/.14}{
  \fill[figureblue,fill opacity=\alpha]
    plot[domain=0:180,samples=73]
      ({2.5*\r*cos(\x)},{3.2*\r+.65*\r*sin(\x)})--
    plot[domain=180:0,samples=73]
      ({2.5*\r*cos(\x)},{3.2*\r*\r+.65*\r*sin(\x)})--cycle;
}
% Group each annular surface so the small patches don't leave dark seams.
\begin{scope}[transparency group,opacity=.14]
\foreach \a in {0,5,...,355}{
  \pgfmathsetmacro{\b}{\a+5}
  \fill[figureblue]
    ({2.5*.58*cos(\a)},{3.2*.58*.58+.65*.58*sin(\a)})--
    ({2.5*.62*cos(\a)},{3.2*.62*.62+.65*.62*sin(\a)})--
    ({2.5*.62*cos(\b)},{3.2*.62*.62+.65*.62*sin(\b)})--
    ({2.5*.58*cos(\b)},{3.2*.58*.58+.65*.58*sin(\b)})--cycle;
}
\end{scope}
\foreach \r/\alpha in {.58/.10,.62/.22}{
  \fill[figureblue,fill opacity=\alpha]
    plot[domain=180:360,samples=73]
      ({2.5*\r*cos(\x)},{3.2*\r+.65*\r*sin(\x)})--
    plot[domain=360:180,samples=73]
      ({2.5*\r*cos(\x)},{3.2*\r*\r+.65*\r*sin(\x)})--cycle;
}
\begin{scope}[transparency group,opacity=.23]
\foreach \a in {0,5,...,355}{
  \pgfmathsetmacro{\b}{\a+5}
  \fill[figureblue]
    ({2.5*.58*cos(\a)},{3.2*.58+.65*.58*sin(\a)})--
    ({2.5*.62*cos(\a)},{3.2*.62+.65*.62*sin(\a)})--
    ({2.5*.62*cos(\b)},{3.2*.62+.65*.62*sin(\b)})--
    ({2.5*.58*cos(\b)},{3.2*.58+.65*.58*sin(\b)})--cycle;
}
\end{scope}
% Draw all four full rims: rear halves dashed, front halves solid.
\begin{scope}[draw=figureaccent,line width=.75pt]
\foreach \r in {.58,.62}{
  \draw[densely dashed,line width=.5pt]
    plot[domain=0:180,samples=73]
      ({2.5*\r*cos(\x)},{3.2*\r+.65*\r*sin(\x)});
  \draw[densely dashed,line width=.5pt]
    plot[domain=0:180,samples=37]
      ({2.5*\r*cos(\x)},{3.2*\r*\r+.65*\r*sin(\x)});
  \draw plot[domain=180:360,samples=73]
    ({2.5*\r*cos(\x)},{3.2*\r+.65*\r*sin(\x)});
  \draw plot[domain=180:360,samples=73]
    ({2.5*\r*cos(\x)},{3.2*\r*\r+.65*\r*sin(\x)});
}
\foreach \s in {-1,1}{
  \draw ({\s*1.55},{3.2*.62*.62})--
    ({\s*1.55},{3.2*.62});
  \draw ({\s*1.45},{3.2*.58*.58})--
    ({\s*1.45},{3.2*.58});
}
\end{scope}
% At the right tangent the inner wall has exactly height f(x)-g(x).
\draw[gray,densely dashed] (1.45,{3.2*.58*.58})--(2.95,{3.2*.58*.58})
  (1.45,{3.2*.58})--(2.95,{3.2*.58});
\draw[<->,black,solid] (2.85,1.07648)--(2.85,1.856);
\node[right=3pt] at (2.85,1.46624) {$h=x-x^2$};
% Radius and thickness refer to the same two radial boundaries.
\draw[gray,densely dashed] (1.45,{3.2*.58*.58})--(1.45,-.65)
  (1.55,{3.2*.62*.62})--(1.55,-.65);
\draw[<->] (0,-.5)--(1.45,-.5) node[midway,below=3pt] {$x$};
\draw[->] (1.05,-.5)--(1.45,-.5);
\draw[->] (2.1,-.5)--(1.55,-.5);
\node[below] at (1.8,-.55) {$dx$};
\draw[->,gray] (0,0)--(0,4.1) node[above,text=black] {$y$};
\end{tikzpicture}
